\documentclass[11pt]{article}

\usepackage[margin=1in]{geometry}
\usepackage{amsmath,amssymb,amsthm}
\usepackage{booktabs}
\usepackage{hyperref}
\usepackage{algorithm}
\usepackage{algpseudocode}

\newtheorem{definition}{Definition}
\newtheorem{proposition}{Proposition}
\newtheorem{remark}{Remark}
\newtheorem{lemma}{Lemma}

\DeclareMathOperator{\RMSE}{RMSE}

\title{Curve-Fitting Inverse Design of a Linearly-Tapered\\
Conical Compression Spring}
\author{Derived from the \texttt{ConicalCurveCompressionSpringInverseDesigner} implementation}
\date{}

\begin{document}
\maketitle

\begin{abstract}
We formalize a regression-based inverse-design procedure for a
linearly-tapered (diameter and pitch) helical compression spring that
must reproduce an entire target force--displacement curve, rather than
two points on its linear regime. Because the curve's nonlinear shape —
the progressive stiffening produced by coil-to-coil and floor contact —
is exactly what encodes the taper, no closed-form solve is available;
instead six geometric parameters (wire diameter, free length, start/end
diameters, start/end pitches) are fit by minimizing a composite objective
combining curve-fit error, a soft fatigue-safety-factor penalty, and a
geometric-feasibility penalty. We first derive a closed-form winding-angle
parametrization for a spring whose pitch varies linearly along its axis,
which replaces a per-point root solve in the general contact-aware
simulator with an explicit exponential map, making thousands of candidate
evaluations inside a global optimizer tractable. We then formalize the
incremental, flexibility-weighted contact simulation used to generate a
candidate's force--displacement curve, the composite regression
objective, and the two-stage (global, then locally polished) optimization
strategy used to solve it.
\end{abstract}

\tableofcontents

\section{Problem statement}

\subsection{Design requirements}

The design brief consists of a material, a target fatigue safety factor
$n_f^{\star}$ (entered as a soft objective term, not a hard constraint),
and a target curve of $N$ samples
\begin{equation}
\big\{(\delta_i, F_i)\big\}_{i=1}^{N}, \qquad \delta_1 < \delta_2 < \dots < \delta_N,
\end{equation}
where $\delta_i$ is displacement (travel) from the free length and $F_i$
the corresponding load.

\subsection{Design variables}

The spring's diameter and pitch are assumed to vary linearly along its
axial coordinate $h\in[0,L_0]$ (free length $L_0$ itself unknown):
\begin{equation}
D(h) = D_0 + \frac{D_1 - D_0}{L_0}\,h, \qquad
p(h) = p_0 + \frac{p_1 - p_0}{L_0}\,h,
\label{eq:linear-profiles}
\end{equation}
with $D_0,D_1$ the start/end mean diameters and $p_0,p_1$ the start/end
pitches. Together with the wire diameter $d$ and $L_0$, the design vector
is
\begin{equation}
x = (d,\, L_0,\, D_0,\, D_1,\, p_0,\, p_1) \in \mathbb{R}^6 .
\label{eq:design-vector}
\end{equation}

\begin{remark}
Unlike the two-point (algebraic) designer for the constant-taper case,
here the rate and its evolution under contact are not linear in the
unknowns, so $x$ cannot be solved for directly; it must be fit by
optimization against the whole curve.
\end{remark}

\section{Closed-form winding-angle parametrization for a linear pitch profile}
\label{sec:closed-form}

\subsection{General relation between winding angle and axial position}

For any pitch profile $p(h)$, the winding angle $\theta$ accumulated
while advancing from the fixed end to axial position $h$ satisfies
\begin{equation}
\theta(h) = \int_0^h \frac{2\pi}{p(s)}\,ds .
\label{eq:theta-general}
\end{equation}
In the general (arbitrary $p(\cdot)$) case, Eq.~\eqref{eq:theta-general}
must be inverted numerically — one root solve per discretization point —
which is the dominant cost of the general contact-aware simulator when
evaluated at the resolution a global search requires.

\subsection{Closed form for a linear pitch profile}

\begin{proposition}
Let $p(h) = p_0 + \alpha h$ with $\alpha = (p_1-p_0)/L_0$ the pitch slope.
For $\alpha \neq 0$,
\begin{equation}
\theta(h) = \frac{2\pi}{\alpha}\,\ln\!\frac{p(h)}{p_0},
\label{eq:theta-h}
\end{equation}
which inverts explicitly to
\begin{equation}
h(\theta) = \frac{p_0}{\alpha}\left(\exp\!\Big(\frac{\alpha\theta}{2\pi}\Big) - 1\right).
\label{eq:h-theta}
\end{equation}
For $\alpha \to 0$ (constant pitch), Eqs.~\eqref{eq:theta-h}--\eqref{eq:h-theta}
degenerate to the linear relation
\begin{equation}
\theta(h) = \frac{2\pi h}{p_0}, \qquad h(\theta) = \frac{p_0\,\theta}{2\pi}.
\label{eq:theta-h-linear}
\end{equation}
\end{proposition}

\begin{proof}
Substituting $p(s) = p_0 + \alpha s$ into Eq.~\eqref{eq:theta-general} and
integrating,
\[
\theta(h) = \frac{2\pi}{\alpha}\Big[\ln(p_0+\alpha s)\Big]_0^{h}
= \frac{2\pi}{\alpha}\ln\frac{p_0+\alpha h}{p_0} = \frac{2\pi}{\alpha}\ln\frac{p(h)}{p_0},
\]
which is Eq.~\eqref{eq:theta-h}. Solving for $h$ gives
Eq.~\eqref{eq:h-theta}. Taking $\alpha\to 0$ in
Eq.~\eqref{eq:theta-general} directly, with $p(s)\equiv p_0$, gives
Eq.~\eqref{eq:theta-h-linear}, which is also the limit of
Eqs.~\eqref{eq:theta-h}--\eqref{eq:h-theta} by L'Hôpital's rule.
\end{proof}

\begin{remark}
Equation~\eqref{eq:h-theta} lets the discretization be sampled uniformly
in winding angle $\theta$ — matching the convention of the general
simulator, so per-turn indexing (points separated by one full winding
turn) still lines up — while avoiding a numerical root solve at every
sample. The maximum winding angle is $\theta_{\max} = \theta(L_0)$ from
Eq.~\eqref{eq:theta-h}, and small numerical excursions of
$h(\theta_{\max})$ outside $[0,L_0]$ are clipped back into range.
\end{remark}

\section{Fast contact-aware progressive compression simulation}
\label{sec:fast-sim}

Given the closed-form samples $\{(\theta_j, h_j)\}_{j=1}^{M}$ from
Section~\ref{sec:closed-form} and the corresponding local radii
$r_j = D(h_j)/2$ (via Eq.~\eqref{eq:linear-profiles}), the spring is
compressed in $S$ fixed axial-displacement increments
$\Delta = L_0/S$, tracking, at each step $t$, a cumulative per-point
axial displacement $y_j^{(t)}$ and an activity flag $a_j^{(t)}\in\{0,1\}$
indicating whether point $j$ is still elastically compliant.

\subsection{Coil-to-coil contact}

Let $q$ denote the number of discretization points per full winding turn
(so index $j$ and $j+q$ sample the same angular phase one turn apart).
At step $t$, the actual (post-deformation) axial position of point $j$ is
\begin{equation}
z_j^{(t)} = h_j - y_j^{(t)} ,
\end{equation}
and the axial gap to the point one turn above is
\begin{equation}
\pi_j^{(t)} = \big|z_{j+q}^{(t)} - z_j^{(t)}\big| .
\label{eq:axial-gap}
\end{equation}
With radial offset $\Delta r_j = |r_{j+q} - r_j|$ (nonzero because the
diameter tapers), two coils of wire diameter $d$ can approach only until
their centerlines are $d$ apart in three dimensions. If $\Delta r_j < d$,
the axial component of that minimum separation is, by the Pythagorean
relation $d^2 = \Delta r_j^2 + \pi_j^{\lim\,2}$,
\begin{equation}
\pi_j^{\lim} = \sqrt{\max\!\big(d^2 - \Delta r_j^2,\ 0\big)} .
\label{eq:pz-limit}
\end{equation}
Contact is detected when $\pi_j^{(t)} \le \pi_j^{\lim}$, in which case
every point in the index range $[j,\,j+q]$ — the entire turn between the
two contacting points — is frozen for the remainder of the step:
\begin{equation}
\pi_j^{(t)} \le \pi_j^{\lim} \implies a_k^{(t)} = 0 \ \text{ for all } k\in[j,j+q].
\label{eq:contact-freeze}
\end{equation}

\subsection{Floor contact}

Any point other than the fixed end ($j\ne 0$) whose actual axial position
has reached or passed zero has bottomed out:
\begin{equation}
z_j^{(t)} \le 0 \ (j \ne 0) \implies a_j^{(t)} = 0 .
\label{eq:floor-contact}
\end{equation}

\subsection{Instantaneous stiffness and load increment}

The local torsional flexibility per unit winding angle at point $j$,
from the standard helical-spring torsion relation, is
\begin{equation}
\phi_j = \frac{8\,D(h_j)^3}{G\,d^{4}\,2\pi},
\label{eq:local-flexibility}
\end{equation}
with $G$ the material shear modulus. Weighting by activity and
integrating over winding angle gives the total instantaneous flexibility
\begin{equation}
\Phi^{(t)} = \int_{0}^{\theta_{\max}} \phi(\theta)\,a^{(t)}(\theta)\,d\theta ,
\label{eq:total-flexibility}
\end{equation}
(evaluated by quadrature over the sampled points), and the instantaneous
stiffness is $k^{(t)} = 1/\Phi^{(t)}$ (or $+\infty$ if
$\Phi^{(t)}\approx 0$, i.e.\ every point has gone rigid). The load
increment over this step is
\begin{equation}
F^{(t)} = F^{(t-1)} + k^{(t-1)}\,\Delta ,
\label{eq:load-increment}
\end{equation}
a small-step linearization of the piecewise-changing stiffness, and the
simulation terminates early if $k^{(t)}=\infty$ (the geometry is fully
solid-packed or floored, so no further deflection is possible).

\subsection{Deflection sharing across active points}

Within a step, the applied deflection increment $\Delta$ is distributed
across active points in proportion to their share of the total
flexibility — the series-spring analogy in which a stiffer (less
flexible) point accepts less of an applied displacement:
\begin{equation}
w_j^{(t)} = \frac{\phi_j\,a_j^{(t)}}{\Phi^{(t)}},
\qquad
Y_j^{(t)} = \int_0^{\theta_j} w^{(t)}(\theta)\,d\theta,
\label{eq:deformation-share}
\end{equation}
(the second integral evaluated by cumulative quadrature over the sampled
winding angles), and the cumulative per-point displacement is updated as
\begin{equation}
y_j^{(t+1)} = y_j^{(t)} + Y_j^{(t)}\,\Delta .
\label{eq:cumulative-update}
\end{equation}

\begin{remark}
Equations~\eqref{eq:axial-gap}--\eqref{eq:cumulative-update} reproduce,
point for point, the contact-detection and load-sharing logic of the
general (arbitrary-profile) progressive-compression simulator; the only
substitution is the closed-form sampling of
Section~\ref{sec:closed-form} in place of a per-point numerical inversion
of Eq.~\eqref{eq:theta-general}. The simulation is therefore a faithful,
not merely approximate, model of the linear-taper geometry — its
computational advantage over the general method is restricted to how the
discretization points are generated, not to any change in the physics.
\end{remark}

\section{Regression objective}
\label{sec:objective}

\subsection{Curve-fit error}

Given a candidate $x$ from Eq.~\eqref{eq:design-vector}, run the
simulation of Section~\ref{sec:fast-sim} to obtain a predicted curve
$(\hat\delta_k,\hat F_k)$, resample it at the target's own displacement
points by linear interpolation,
\begin{equation}
\hat F(\delta_i) = \operatorname{interp}\big(\delta_i;\ \{\hat\delta_k\},\{\hat F_k\}\big),
\end{equation}
and compute the root-mean-square error against the target curve,
\begin{equation}
\RMSE(x) = \sqrt{\frac{1}{N}\sum_{i=1}^N \big(\hat F(\delta_i) - F_i\big)^2}.
\label{eq:rmse}
\end{equation}
Normalizing by the target curve's typical magnitude,
\begin{equation}
\bar F = \max\!\left(\frac{1}{N}\sum_{i=1}^N |F_i|,\ \varepsilon\right)
\qquad (\varepsilon = 10^{-6}),
\end{equation}
gives a unitless relative error
\begin{equation}
\widetilde{\RMSE}(x) = \frac{\RMSE(x)}{\bar F},
\label{eq:norm-rmse}
\end{equation}
comparable across curves of very different force ranges.

\subsection{Soft fatigue penalty}

At the two extremes of the target displacement range, the critical-section
stresses are evaluated using the local diameter of
Eq.~\eqref{eq:linear-profiles} and the simulated load there, via the
Wahl-corrected torsion formula (as in the companion designer papers),
yielding a safety factor $n_f(x)$ from the Goodman criterion. The
deviation from target enters the objective as a soft penalty,
\begin{equation}
\Pi_{\mathrm{fatigue}}(x) = \big|n_f(x) - n_f^{\star}\big| .
\label{eq:fatigue-penalty}
\end{equation}

\subsection{Geometric feasibility penalty}

Let $C_{\min},C_{\max}$ bound the admissible spring index and $n^{\min}$
the minimum coil count. With coil count
\begin{equation}
n(x) = \frac{L_0}{p_0}\int_0^1 \frac{p_0}{p(u L_0)/p_0}\,du
\end{equation}
(the linear-taper analogue of active coils, from integrating winding
turns along the taper) and solid length
\begin{equation}
L_s(x) =
\begin{cases}
d, & |D_0-D_1| \ge n(x)\,d \quad \text{(telescoping)}, \\
n(x)\,d, & \text{otherwise (stacking)},
\end{cases}
\label{eq:solid-length-curve}
\end{equation}
(the same criterion as Eq.~\eqref{eq:solid-length-inv} of the point-based
designer, applied at both taper ends), the feasibility penalty sums four
smoothly-vanishing (squared, one-sided) violation terms:
\begin{align}
\Pi_{\mathrm{geom}}(x) = {}& \sum_{D\in\{D_0,D_1\}}\Big[\max(0, C_{\min}-D/d)^2 + \max(0, D/d-C_{\max})^2\Big] \notag\\
&+ \sum_{p\in\{p_0,p_1\}} \max(0, d-p)^2 \notag \\
&+ \max\!\big(0,\ n^{\min}-n(x)\big)^2
+ \max\!\big(0,\ L_s(x) - \ell_{\mathrm{hi}}\big)^2 ,
\label{eq:geometry-penalty}
\end{align}
where $\ell_{\mathrm{hi}} = L_0 - \delta_{\min}$ is the shortest length
the target curve exercises (from the smallest target displacement
$\delta_{\min}$). Each term vanishes at its constraint boundary and grows
quadratically beyond it, so the penalty is continuous and
"differentiable-ish" — supporting the local polish stage of
Section~\ref{sec:optimization} — rather than a hard reject.

\subsection{Composite objective}

The full objective minimized by the search is
\begin{equation}
J(x) = \widetilde{\RMSE}(x) + w_f\,\Pi_{\mathrm{fatigue}}(x) + w_g\,\Pi_{\mathrm{geom}}(x),
\label{eq:composite-objective}
\end{equation}
with weights $w_f,w_g>0$, and any evaluation failure (a degenerate or
numerically invalid candidate geometry) is assigned a large constant
value $J=10^{6}$ so the optimizer treats it as strictly worse than any
successful evaluation without terminating the search.

\section{Two-stage optimization}
\label{sec:optimization}

\subsection{Search bounds}

Each of the six design variables is bounded: wire diameter and taper
diameters/pitches within user-supplied ranges, and free length within a
margin of the target curve's maximum displacement $\delta_{\max}$,
\begin{equation}
L_0 \in \big[m_{\mathrm{lo}}\,\delta_{\max},\ m_{\mathrm{hi}}\,\delta_{\max}\big], \qquad m_{\mathrm{lo}} > 1,
\label{eq:free-length-bounds}
\end{equation}
the lower margin exceeding unity because the spring cannot be compressed
past its own free length.

\subsection{Global stage}

The composite objective $J(x)$ (Eq.~\eqref{eq:composite-objective}),
evaluated with a coarse simulation resolution
(Section~\ref{sec:fast-sim} at reduced point/step counts), is minimized
over the bounded box by a derivative-free global optimizer (differential
evolution),
\begin{equation}
x^{(0)} = \operatorname*{arg\,min}_{x \in \mathcal{B}} J(x),
\label{eq:global-stage}
\end{equation}
$\mathcal{B}\subset\mathbb{R}^6$ the bound box of
Section~\ref{sec:objective}--\ref{eq:free-length-bounds}. This handles
the non-convex, discontinuous (due to the contact-freezing logic)
landscape of $J$ without requiring gradients.

\subsection{Diameter snapping and local polish}

The continuous wire diameter $d^{(0)}$ from $x^{(0)}$ is snapped to the
nearest manufacturable size in a standard series $\mathcal{D}_{\mathrm{std}}$,
\begin{equation}
d^{\star} = \operatorname*{arg\,min}_{d \in \mathcal{D}_{\mathrm{std}}} \big|d - d^{(0)}\big| .
\label{eq:diameter-snap}
\end{equation}
Holding $d=d^{\star}$ fixed, the remaining five parameters
$y = (L_0,D_0,D_1,p_0,p_1)$ are then re-optimized locally by a
derivative-free simplex method (Nelder--Mead), started from $x^{(0)}$'s
values, to recover fit quality lost by the diameter snap:
\begin{equation}
y^{\star} = \operatorname*{arg\,min}_{y} J\big(d^{\star}, y\big).
\label{eq:local-polish}
\end{equation}

\subsection{Final verification}

The optimized geometry $(d^{\star}, y^{\star})$ is rebuilt using the
general, contact-aware spring model at full (fine) simulation resolution
— not the fast closed-form approximation of Section~\ref{sec:fast-sim} —
and its curve, stresses, and Goodman safety factor are recomputed from
that unmodified machinery, at every target displacement point (not just
the two curve extremes), so that everything reported to the caller is
verified rather than approximated.

\begin{algorithm}[h]
\caption{Curve-fitting inverse design of a linear-taper conical spring}
\begin{algorithmic}[1]
\Require material, target curve $\{(\delta_i,F_i)\}$, $n_f^{\star}$, bounds
\State Normalize $\bar F$ from the target curve; set search resolution (coarse) and final resolution (fine)
\State $x^{(0)} \gets$ global minimization of $J$ (Eq.~\eqref{eq:composite-objective}) via differential evolution over $\mathcal{B}$ (coarse simulation, Section~\ref{sec:fast-sim})
\State $d^{\star} \gets$ nearest standard wire diameter to $x^{(0)}$'s wire diameter (Eq.~\eqref{eq:diameter-snap})
\State $y^{\star} \gets$ local Nelder--Mead re-fit of the remaining five parameters at $d=d^{\star}$ (Eq.~\eqref{eq:local-polish})
\State Build the general contact-aware spring at $(d^{\star},y^{\star})$; simulate at fine resolution
\State Register every target displacement as a load position on the real spring; recompute safety factor from its own stress machinery
\State \Return spring, geometry, curve fit ($\RMSE$, $\widetilde{\RMSE}$), safety factor, both curves
\end{algorithmic}
\end{algorithm}

\section{Discussion}

\begin{proposition}
The closed-form map of Eqs.~\eqref{eq:theta-h}--\eqref{eq:h-theta}
reduces the per-evaluation cost of the contact simulation from one
root-finding call per discretization point to a single vectorized
exponential evaluation, without altering the contact-detection or
load-sharing physics of Section~\ref{sec:fast-sim}, which are identical
to the general (arbitrary-profile) simulator. It is this reduction that
makes a several-thousand-evaluation global search
(Eq.~\eqref{eq:global-stage}) computationally tractable at all, since the
general per-point-fsolve simulator would otherwise dominate the total
search time by orders of magnitude.
\end{proposition}

\begin{remark}
The design problem is over-determined in the sense that a full curve —
not merely a rate and a free length — must be matched by only six
parameters under a linear-taper assumption; consequently the fitted
curve should be expected to approximate, rather than exactly reproduce,
target curves that are not themselves the output of a linear-taper
geometry. The quantities $\RMSE(x^{\star})$ and
$\widetilde{\RMSE}(x^{\star})$ of Eq.~\eqref{eq:rmse}--\eqref{eq:norm-rmse},
evaluated on the final verified curve, report the residual fit quality
for exactly this reason.
\end{remark}

\section{Conclusion}

The inverse curve-fitting design problem for a linearly-tapered
compression spring is solved by (i) an exact reparametrization of the
winding angle for linear pitch profiles
(Eqs.~\eqref{eq:theta-h}--\eqref{eq:h-theta}), which makes a
high-fidelity contact simulation cheap enough to evaluate thousands of
times; (ii) a discrete-step, flexibility-weighted simulation of
progressive coil-to-coil and floor contact
(Eqs.~\eqref{eq:axial-gap}--\eqref{eq:cumulative-update}) that produces
each candidate's force--displacement curve; (iii) a composite objective
blending normalized curve-fit error, a soft fatigue-safety-factor term,
and a smooth geometric-feasibility penalty
(Eq.~\eqref{eq:composite-objective}); and (iv) a two-stage
global-then-local optimization (Eqs.~\eqref{eq:global-stage}--\eqref{eq:local-polish})
that respects a discrete manufacturable wire-diameter series while
recovering as much fit quality as the continuous parameters allow.

\end{document}
